\documentclass[11pt,a4paper]{article}

\usepackage[utf8]{inputenc}
\usepackage[T1]{fontenc}
\usepackage[french]{babel}
\usepackage[a4paper,margin=2cm]{geometry}
\usepackage{amsmath,amssymb}
\usepackage{lmodern}
\usepackage{enumitem}
\usepackage{array}
\usepackage{tabularx}
\usepackage{xcolor}
\usepackage{fancyhdr}
\usepackage{lastpage}
\usepackage{tikz}
\usepackage{multicol}
\usepackage[hidelinks]{hyperref}

\setlength{\parindent}{0pt}
\setlength{\parskip}{0.4em}

\pagestyle{fancy}
\fancyhf{}
\lhead{DS de mathématiques -- Niveau 3\ieme}
\rhead{Nom : \hspace{4cm}}
\cfoot{\thepage/\pageref{LastPage}}

\newcommand{\ds}{\displaystyle}

\begin{document}

\begin{center}
    {\Large \textbf{DS de mathématiques}}\\[0.3em]
    {\large \textbf{Niveau 3\ieme}}\\[0.5em]
\end{center}

\begin{center}
\renewcommand{\arraystretch}{1.4}
\begin{tabular}{|>{\centering\arraybackslash}m{3cm}|>{\centering\arraybackslash}m{3cm}|>{\centering\arraybackslash}m{6cm}|}
\hline
\textbf{Durée} & \textbf{Barème} & \textbf{Matériel autorisé} \\
\hline
1 h 30 à 2 h & 20 points & Stylo, règle, équerre, compas, calculatrice \\
\hline
\end{tabular}
\end{center}

\textbf{Consignes générales}
\begin{itemize}[leftmargin=1.2cm]
    \item Pensez bien à noter le temps que vous avez pris pour faire le sujet.
    \item Toute réponse doit être justifiée.
    \item Les calculs doivent être rédigés clairement.
    \item Lorsqu’un rappel ou une aide est donné, il peut être utilisé librement.
    \item La qualité de la rédaction et le soin apporté à la copie seront pris en compte.
\end{itemize}

\vspace{0.3cm}

\begin{center}
\renewcommand{\arraystretch}{1.4}
\begin{tabular}{|>{\centering\arraybackslash}m{2.2cm}|>{\centering\arraybackslash}m{2.2cm}|>{\centering\arraybackslash}m{2.2cm}|>{\centering\arraybackslash}m{2.2cm}|>{\centering\arraybackslash}m{2.2cm}|>{\centering\arraybackslash}m{2.2cm}|}
\hline
\textbf{Exercice 1} & \textbf{Exercice 2} & \textbf{Exercice 3} & \textbf{Exercice 4} & \textbf{Exercice 5} & \textbf{Total} \\
\hline
5 pts & 4 pts & 4 pts & 3 pts & 4 pts & 20 pts \\
\hline
\end{tabular}
\end{center}

\section*{Exercice 1 -- Calcul littéral, équation et fonction \hfill /5}

On considère la fonction \(f\) définie par :
\[
f(x)=x^2-5x+6.
\]

\textbf{Rappels utiles}
\begin{itemize}[leftmargin=1.2cm]
    \item Calculer l’image d’un nombre par une fonction, c’est remplacer \(x\) par ce nombre.
    \item Si un produit est nul, alors au moins un des facteurs est nul.
    \item Une expression développée et une expression factorisée peuvent représenter la même fonction.
\end{itemize}

\begin{enumerate}
    \item Calculer \(f(0)\), \(f(1)\), \(f(2)\) et \(f(3)\).
    \item Montrer que, pour tout réel \(x\),
    \[
    f(x)=(x-2)(x-3).
    \]
    \item Résoudre l’équation
    \[
    f(x)=0.
    \]
    \item Résoudre l’équation
    \[
    f(x)=6.
    \]
    \item Déterminer les réels \(x\) tels que
    \[
    f(x)<0.
    \]
\end{enumerate}

\textbf{Barème indicatif :} 1 pt ; 1 pt ; 1 pt ; 1 pt ; 1 pt.

\section*{Exercice 2 -- Géométrie : théorème de Thalès \hfill /4}

Dans la figure ci-dessous, les points \(A\), \(B\), \(M\) sont alignés, les points \(A\), \(C\), \(N\) sont alignés et les droites \((BC)\) et \((MN)\) sont parallèles.

On donne :
\[
AB=6 \text{ cm}, \qquad AM=9 \text{ cm}, \qquad AC=4 \text{ cm}.
\]

\textbf{Rappels utiles}
\begin{itemize}[leftmargin=1.2cm]
    \item Si \(B\in[AM]\), \(C\in[AN]\) et \((BC)\parallel(MN)\), alors
    \[
    \frac{AB}{AM}=\frac{AC}{AN}=\frac{BC}{MN}.
    \]
    \item On pense à bien vérifier les conditions d’utilisation du théorème de Thalès.
\end{itemize}

\begin{enumerate}
    \item Expliquer pourquoi on peut utiliser le théorème de Thalès dans cette situation.
    \item Calculer la longueur \(AN\).
    \item On donne de plus \(BC=5\) cm. Calculer la longueur \(MN\).
    \item Sans refaire tous les calculs, expliquer si le triangle \(AMN\) est un agrandissement ou une réduction du triangle \(ABC\).
\end{enumerate}

\textbf{Barème indicatif :} 1 pt ; 1,5 pt ; 1 pt ; 0,5 pt.

\section*{Exercice 3 -- Démonstration : théorème de Pythagore \hfill /4}

On considère un triangle \(ABC\) rectangle en \(A\), tel que :
\[
AB=6 \text{ cm} \qquad \text{et} \qquad AC=8 \text{ cm}.
\]

\textbf{Rappels utiles}
\begin{itemize}[leftmargin=1.2cm]
    \item Dans un triangle rectangle, le théorème de Pythagore dit que :
    \[
    \text{(hypoténuse)}^2=\text{(côté 1)}^2+\text{(côté 2)}^2.
    \]
    \item La réciproque permet de montrer qu’un triangle est rectangle.
\end{itemize}

\begin{enumerate}
    \item Calculer la longueur \(BC\).
    \item Le point \(D\) est tel que \(AD=2\) cm et \(D\in[AB]\). Faire un schéma soigné.
    \item Calculer la longueur \(DC\).
    \item Le triangle \(ADC\) est-il rectangle ? Justifier.
\end{enumerate}

\textbf{Barème indicatif :} 1,5 pt ; 0,5 pt ; 1 pt ; 1 pt.

\section*{Exercice 4 -- Statistiques et probabilités \hfill /3}

Une classe de 25 élèves a obtenu les notes suivantes à un contrôle :

\begin{center}
\begin{tabular}{|c|c|c|c|c|c|c|}
\hline
Note & 6 & 8 & 10 & 12 & 15 & 18 \\
\hline
Effectif & 2 & 4 & 7 & 6 & 4 & 2 \\
\hline
\end{tabular}
\end{center}

\textbf{Rappels utiles}
\begin{itemize}[leftmargin=1.2cm]
    \item La moyenne est la somme des notes pondérées par leurs effectifs, divisée par l’effectif total.
    \item La médiane partage la série en deux groupes de même effectif.
    \item Une probabilité se calcule par :
    \[
    \frac{\text{nombre de cas favorables}}{\text{nombre de cas possibles}}
    \]
    lorsque les issues sont équiprobables.
\end{itemize}

\begin{enumerate}
    \item Vérifier que l’effectif total est bien 25.
    \item Calculer la moyenne de la classe.
    \item On choisit au hasard un élève de la classe. Quelle est la probabilité qu’il ait obtenu au moins 12 ?
\end{enumerate}

\textbf{Barème indicatif :} 0,5 pt ; 1,5 pt ; 1 pt.

\section*{Exercice 5 -- Logique, nombres et raisonnement \hfill /4}

\textbf{Rappels utiles}
\begin{itemize}[leftmargin=1.2cm]
    \item Un nombre pair s’écrit sous la forme \(2n\), avec \(n\in\mathbb{Z}\).
    \item Un nombre impair s’écrit sous la forme \(2n+1\), avec \(n\in\mathbb{Z}\).
    \item Si un produit est nul, alors au moins un des facteurs est nul.
\end{itemize}

\begin{enumerate}
    \item Montrer que la somme de deux nombres impairs est paire.
    \item Montrer que le carré d’un nombre pair est pair.
    \item Résoudre l’équation
    \[
    (x-4)(2x+3)=0.
    \]
    \item On considère l’affirmation suivante :\\
    \emph{« Si un nombre est multiple de 6, alors il est multiple de 3. »}\\
    Dire si cette affirmation est vraie ou fausse, puis justifier.
\end{enumerate}

\textbf{Barème indicatif :} 1 pt ; 1 pt ; 1 pt ; 1 pt.

\vspace{0.5cm}

\begin{center}
\textbf{Fin du sujet -- Toute réponse doit être justifiée avec soin.}
\end{center}

\end{document}